% GENERATED FILE. Edit canonical lessons, then run npm run generate:books.
% canonical-source-hash: fnv1a64:c174deeb8c7e699f
% canonical-lessons: FR-C289-agression, FR-C289-agresser, FR-C289-agresseur, FR-C289-menacer, FR-C289-menace

\chapter{Une agression}
\label{ch:fr-une-agression}
\hlchaptermodality{289}

\begin{chapteropening}
\textbf{By the end of this chapter:} \emph{I can report an attack: the assault, attacking, the attacker, threatening and a threat.}

Complete the last lesson of chapter 289: I can report an attack: the assault, attacking, the attacker, threatening and a threat.
\end{chapteropening}

\section[l'agression]{l'agression — the attack, the assault}

\label{lesson:FR-C289-agression}



\begin{quote}
Before the new one: say the French for to swindle, then the French for the fraud.
\end{quote}

\subsection*{You'll want to know: l'agression}
\textbf{l'agression} — \textquotedblleft{}the attack, the assault\textquotedblright{}.

\begin{quote}
Une agression a eu lieu près du marché.
\end{quote}

\begin{quote}
\emph{An assault took place near the market.}
\end{quote}

An attack on a person. A hostile mood is \emph{l'agressivité}.

\begin{cousinweb}
\emph{agression} is a learned borrowing from Latin \emph{aggressiō} \textquotedblleft{}a going towards\textquotedblright{}, from \emph{aggredī} \textquotedblleft{}to approach, to attack\textquotedblright{}, built on \emph{gradī} \textquotedblleft{}to step\textquotedblright{} (Lewis \& Short).
\end{cousinweb}

\begin{grammarlens}[title={What you've built}]
One more word for telling, explaining and describing at B1.
\end{grammarlens}

\subsection*{Your turn}
\begin{itemize}
\raggedright
  \item \emph{Say it:} \emph{l'agression}
  \item \emph{Say it:} \emph{l'agression}, once more
  \item \emph{Say it:} a sentence of your own with \emph{l'agression}
  \item \emph{Recall:} say the French for to swindle, then the French for the fraud, then say \emph{l'agression} again
\end{itemize}

\begin{tcolorbox}[breakable,colback=teal!4,colframe=teal!35!black,title={Before you move on}]
What does l'agression mean? (\textquotedblleft{}The attack, the assault\textquotedblright{}.) Say it once more.
\end{tcolorbox}
\section[agresser]{agresser — to attack, to assault}

\label{lesson:FR-C289-agresser}



\begin{quote}
Before the new one: say the French for the fraud, then the French for the attack.
\end{quote}

\subsection*{You'll want to know: agresser}
\textbf{agresser} — \textquotedblleft{}to attack, to assault\textquotedblright{}.

\begin{quote}
Deux hommes ont agressé un jeune homme dans la rue.
\end{quote}

\begin{quote}
\emph{Two men attacked a young man in the street.}
\end{quote}

Also figuratively: \emph{ce bruit m'agresse} — this noise gets to me.

\begin{cousinweb}
\emph{agresser} is formed on the stem of \emph{agression}, a learned borrowing from Latin \emph{aggressiō}, from \emph{aggredī} \textquotedblleft{}to attack\textquotedblright{} (Lewis \& Short).
\end{cousinweb}

\begin{grammarlens}[title={What you've built}]
One more word for telling, explaining and describing at B1.
\end{grammarlens}

\subsection*{Your turn}
\begin{itemize}
\raggedright
  \item \emph{Say it:} \emph{agresser}
  \item \emph{Say it:} \emph{agresser}, once more
  \item \emph{Say it:} a sentence of your own with \emph{agresser}
  \item \emph{Recall:} say the French for the fraud, then the French for the attack, then say \emph{agresser} again
\end{itemize}

\begin{tcolorbox}[breakable,colback=teal!4,colframe=teal!35!black,title={Before you move on}]
What does agresser mean? (\textquotedblleft{}To attack, to assault\textquotedblright{}.) Say it once more.
\end{tcolorbox}
\section[l'agresseur]{l'agresseur — the attacker, the assailant}

\label{lesson:FR-C289-agresseur}



\begin{quote}
Before the new one: say the French for the attack, then the French for to attack.
\end{quote}

\subsection*{You'll want to know: l'agresseur}
\textbf{l'agresseur} — \textquotedblleft{}the attacker, the assailant\textquotedblright{}.

\begin{quote}
L'agresseur est parti en courant.
\end{quote}

\begin{quote}
\emph{The attacker ran off.}
\end{quote}

In news reports the attacker is often unnamed: \emph{un agresseur inconnu}.

\begin{cousinweb}
\emph{agresseur} is a learned borrowing from Latin \emph{aggressor} \textquotedblleft{}an attacker\textquotedblright{}, from \emph{aggredī} \textquotedblleft{}to attack\textquotedblright{} (Lewis \& Short).
\end{cousinweb}

\begin{grammarlens}[title={What you've built}]
One more word for telling, explaining and describing at B1.
\end{grammarlens}

\subsection*{Your turn}
\begin{itemize}
\raggedright
  \item \emph{Say it:} \emph{l'agresseur}
  \item \emph{Say it:} \emph{l'agresseur}, once more
  \item \emph{Say it:} a sentence of your own with \emph{l'agresseur}
  \item \emph{Recall:} say the French for the attack, then the French for to attack, then say \emph{l'agresseur} again
\end{itemize}

\begin{tcolorbox}[breakable,colback=teal!4,colframe=teal!35!black,title={Before you move on}]
What does l'agresseur mean? (\textquotedblleft{}The attacker, the assailant\textquotedblright{}.) Say it once more.
\end{tcolorbox}
\section[menacer]{menacer — to threaten}

\label{lesson:FR-C289-menacer}



\begin{quote}
Before the new one: say the French for to attack, then the French for the attacker.
\end{quote}

\subsection*{You'll want to know: menacer}
\textbf{menacer} — \textquotedblleft{}to threaten\textquotedblright{}.

\begin{quote}
Il a menacé le vendeur avec un couteau.
\end{quote}

\begin{quote}
\emph{He threatened the shop assistant with a knife.}
\end{quote}

\emph{Menacer de} + infinitive: \emph{il a menacé de partir}.

\begin{cousinweb}
\emph{menacer} is inherited from Vulgar Latin \emph{mināciāre} \textquotedblleft{}to threaten\textquotedblright{} (EtymDB), built on Latin \emph{mināx} \textquotedblleft{}threatening\textquotedblright{}.
\end{cousinweb}

\begin{grammarlens}[title={What you've built}]
One more word for telling, explaining and describing at B1.
\end{grammarlens}

\subsection*{Your turn}
\begin{itemize}
\raggedright
  \item \emph{Say it:} \emph{menacer}
  \item \emph{Say it:} \emph{menacer}, once more
  \item \emph{Say it:} a sentence of your own with \emph{menacer}
  \item \emph{Recall:} say the French for to attack, then the French for the attacker, then say \emph{menacer} again
\end{itemize}

\begin{tcolorbox}[breakable,colback=teal!4,colframe=teal!35!black,title={Before you move on}]
What does menacer mean? (\textquotedblleft{}To threaten\textquotedblright{}.) Say it once more.
\end{tcolorbox}
\section[la menace]{la menace — the threat}

\label{lesson:FR-C289-menace}



\begin{quote}
Before the new one: say the French for the attacker, then the French for to threaten.
\end{quote}

\subsection*{You'll want to know: la menace}
\textbf{la menace} — \textquotedblleft{}the threat\textquotedblright{}.

\begin{quote}
Elle a reçu des menaces au téléphone.
\end{quote}

\begin{quote}
\emph{She got threats over the phone.}
\end{quote}

\emph{Une menace de mort} — a death threat.

\begin{cousinweb}
\emph{menace} is inherited from Vulgar Latin \emph{minācia} \textquotedblleft{}a threat\textquotedblright{} (EtymDB), built on Latin \emph{mināx} \textquotedblleft{}threatening\textquotedblright{}.
\end{cousinweb}

\begin{grammarlens}[title={What you've built}]
One more word for telling, explaining and describing at B1.
\end{grammarlens}

\subsection*{Your turn}
\begin{itemize}
\raggedright
  \item \emph{Say it:} \emph{la menace}
  \item \emph{Say it:} \emph{la menace}, once more
  \item \emph{Say it:} a sentence of your own with \emph{la menace}
  \item \emph{Recall:} say the French for the attacker, then the French for to threaten, then say \emph{la menace} again
\end{itemize}

\begin{tcolorbox}[breakable,colback=teal!4,colframe=teal!35!black,title={Before you move on}]
What does la menace mean? (\textquotedblleft{}The threat\textquotedblright{}.) Say it once more.
\end{tcolorbox}
